\section{覆盖与可区分性：规则~110 中的构造性见证}
\label{sec:observer-witnesses}

同一动力系统可以提出两类问题。在第一类问题中，观察者根据若干格点的读数记录，预测某个元胞未来的值。在第二类问题中，局部记录已无法区分一个过去比特 $Y$ 原先的两个可能取值，接收方却需要恢复自身的这一比特。两类问题都要检验第~\ref{sec:observer-theory} 节的判据：目标在观测映射的每个纤维上是否取常值。

\subsection{根据记录作预测，以及几何覆盖的局限}
\label{sec:observer-coverage}

对于初始构型 $b$，用 $b_t(s)$ 表示时刻 $t$ 位置 $s$ 处的值。设 $S$ 为有限的固定位置列表，$\tau$ 为最后一次读取的时刻。观测记录为数组
\[
 H_{S,\tau}(b)=\bigl(b_k(i)\bigr)_{0\le k\le\tau,\,i\in S}.
\]
在源文件 \nolinkurl{observer/Rule110Observer.lean} 中，谓词 \texttt{CanPredict} 要求存在一个解码器，使其能在所有允许构型上由此记录给出 $b_t(s)$。定理 \texttt{canPredict\_iff} 表明，这一要求等价于 $b_t(s)$ 在 $H_{S,\tau}$ 的每个纤维上取常值。定义允许对初始构型施加任意限制。

用 $I_t(s)=s+I_t$ 表示平移后的精确依赖区域。我们使用确立 $D_t=I_t$ 的定理 \texttt{Rails.essential\_iff}，并利用动力系统的平移不变性将结论从原点移至 $s$。谓词 \texttt{CoversTarget} 的含义是
\begin{equation}
 I_t(s)\ \subseteq\ \bigcup_{i\in S}\bigcup_{k=0}^{\tau} I_k(i).
 \label{eq:observer-cover}
\end{equation}
对于全体初始构型，这是一个必要条件。若左侧的位置 $j$ 不属于右侧，则由 $j$ 的本质性可知，存在两个仅在 $j$ 处不同、但 $b_t(s)$ 不同的构型。它们的全部读数相同，因为 $j$ 位于每次读取的依赖区域之外。这对构型违反了可预测性判据。允许采用哪个初始状态类在此至关重要：若该类只有一个构型，则任一未来比特都可用一个常量预测。

\paragraph{覆盖不充分的反例。}
取 $S=(-1,1)$、$\tau=1$、$s=0$、$t=2$。仅最后两处读取的依赖锥 $[-2,0]$ 与 $[0,2]$ 就已覆盖 $I_2(0)=[-2,2]$。第一个初始构型只有位置 $-1$ 处为 1，第二个构型在 $-1$ 和 $0$ 处为 1。在这两种情况下，其他位置均为零。经过一次更新，取值为 1 的位置集合分别为 $\{-2,-1\}$ 和 $\{-2,-1,0\}$。两个世界中的传感器都产生同一份记录 $(10,10)$。在位置 $0$ 的第二次更新中，两个局部三元组 $100$ 和 $110$ 分别给出 $0$ 和 $1$。由此可见，尽管覆盖条件成立，依据完整记录仍无法预测目标。精确构造见 \texttt{separated\_cells\_counterexample} 和 \texttt{coverage\_does\_not\_suffice}。

采用相邻的传感器也会出现这一现象。取 $S=(0,1)$ 和 $\tau=2$，则依赖锥 $[-2,2]$ 与 $[-1,3]$ 覆盖 $I_3(0)=[-2,3]$。初始取值为 1 的位置集合分别为 $\{2\}$ 和 $\{-1,2\}$ 时，所得记录都是 $(00,01,11)$，而时刻 $3$ 位置 $0$ 的取值分别为 $1$ 和 $0$。这就是 \texttt{adjacent\_cells\_counterexample}。也有一个意义直接的充分条件：若在最后一次观测时读出了 $I_{t-\tau}(s)$ 中\emph{所有元胞}的值，且 $t\ge\tau$，则由这些值出发的正常演化唯一确定目标。式~\eqref{eq:observer-cover} 允许将初始数据压缩为少数元胞读数，而这一压缩过程可能使原有的差异无法区分。

\subsection{四段成功的历史与完整的 11 比特对照}

在整数直线上，将一个 11 比特字从左到右置于位置 $0,\ldots,10$，字以外的背景为零。在第一次更新前，把位置 $11$ 置为 1。过去的目标比特为 $Y=b_0(5)$，传感器读数为 $Z=b_2(10)$，接收方记录为 $H=(b_k(0))_{k=0}^3$，记录的最后一个比特为 $K=b_3(0)$。在时刻 3 的观测之后，把 $Z$ 写入 $-1$，并读取该元胞在时刻 4 经常规更新得到的值。

\begin{center}
\begin{tabular}{lccccc}
初始字 & $H$ & $K$ & $Y$ & $Z$ & 时刻 4 的输出 \\
\hline
\texttt{10111010001} & 1100 & 0 & 0 & 0 & 0 \\
\texttt{10111110111} & 1100 & 0 & 1 & 1 & 1 \\
\texttt{11000010111} & 1101 & 1 & 0 & 1 & 0 \\
\texttt{11001100001} & 1101 & 1 & 1 & 0 & 1
\end{tabular}
\end{center}

这四段历史均满足 $Y=K\oplus Z$。对每个固定的 $K$，在时刻 3 执行写入操作之前，一对历史中甚至整个主构型都已相同。接收元胞的左邻居为 1，右邻居为 $K$。因此，写入的比特会依据下述恒等式产生正确答案：
\begin{equation}
 f_{110}(1,z,k)=z\oplus k.
 \label{eq:observer-xor-gate}
\end{equation}
在这四种情况下，将送达的比特取反都会使答案翻转。不过，每一对预先构造的字在位置 $5$ 之外也有差异。

在由全部 $2^{11}$ 个字构成的均匀样本空间上，将每个字与仅翻转 $b_0(5)$ 所得的副本进行比较。时刻 2 位置 $10$ 的传感器依赖于初始位置 $8,\ldots,12$，而位置 $0$ 截至时刻 3 的记录完全由位置 $0$ 及其三步邻域确定。因此，在全部 1024 对中，每对里的两个 $Y$ 值都产生相同的观测 $(H,Z)$。任何基于 $(H,Z)$ 的确定性二元回答，都恰好在 1024/2048 段历史上正确。\nolinkurl{observer/domain.py} 的穷举确认，读数改变的情形为 0/1024，解码器的正确率上限为 1024/2048。这区分了两个命题：\emph{写入的}比特对接收元胞的因果效应，以及\emph{过去的}比特对传感器的因果效应。

\subsection{单独翻转初始比特并在时刻 5 恢复目标}
\label{sec:observer-return}

现在将下表中的字置于从 $-6$ 到 $13$ 的位置上（含两端），区间以外保持零背景。在此构造中，观察者被动读取 $Z=b_2(3)$。在比特送回之前，$Y=b_0(5)$，$K=b_4(0)$。对于每个固定的 $K$，两个字仅在位置 $5$ 不同。

\begin{center}
\begin{tabular}{cc l cc}
$K$ & $Y$ & $[-6,13]$ 上的字 & $Z$ & 时刻 5 的输出 \\
\hline
0 & 0 & \texttt{00000001100001000000} & 0 & 0 \\
0 & 1 & \texttt{00000001100101000000} & 1 & 1 \\
1 & 0 & \texttt{00010111010001000010} & 1 & 0 \\
1 & 1 & \texttt{00010111010101000010} & 0 & 1
\end{tabular}
\end{center}

对于 $K$ 的两个取值，单独翻转 $Y$ 都会改变 $Z$。当 $K=0$ 时，传感器返回 $Y$；当 $K=1$ 时，返回它的补值。到时刻 4，每对历史中的主构型已在整条直线上一致。写入之前，从时刻 0 到 4，接收窗口 $(-2,-1,0)$ 的记录也完全一致。此时 $b_4(-2)=1$。将 $Z$ 一次性写入位置 $-1$，再按规则~110 更新一步，由式~\eqref{eq:observer-xor-gate} 得 $b_5(-1)=Z\oplus K=Y$。若写入 $1-Z$，则四个世界都得到 $1-Y$。写入之前，主层已经丢失上述差异，外部的单比特寄存器则保留了它。无限直线上的计算由 \nolinkurl{observer/causal_line.py} 实现。

\paragraph{最小性结论的适用范围。}
对于直线上满足以下条件的一类构造，时刻 4 是完全合并的最早时刻：每对历史仅翻转一个初始元胞，在预先固定的位置和时刻取得一次原始读数，而且要在四段历史构成的方阵上精确恢复目标，两个二元输入 $K,Z$ 都不可缺少。考察到时刻 3 是否已经合并，只需看该时刻翻转点周围半径为 3 的区间，因为所有可能的差异均位于其中，并由初始时半径为 6 的区间决定。去掉被翻转的元胞，还剩 12 个环境比特。在 $2^{12}=4096$ 种环境中，有八种在第三步或此前就实现完全合并。选定一个敏感的传感器，并固定合并前的读数时刻。在上述环境中，凡是该传感器能感知翻转，其响应方向就都相同。为精确作答，传感器在两个上下文中都必须敏感。此时若响应方向相同，单凭 $Z$ 就能确定 $Y$，与两个输入均不可缺少的条件相矛盾。表中的四个字达到了合并时间的下界，并在下一步更新时给出输出。环境的穷举见 \nolinkurl{observer/causal_crosscheck.py}。

\subsection{有限速度传递与自主执行}

在上一个实验中，外部装置将已保存的比特直接写入接收位置。对相同的四个初始字，还可以显式限制信道的传递速度。在时刻 $2,3,4,5,6,7$，被隔离的比特依次位于 $3,3,3,2,1,0$。所有世界采用相同的时间安排，每步最多移动一个元胞。时刻 7，元胞 $0$ 两侧的值为 $b_7(-1)=1$ 和 $b_7(1)=K$。写入送达的 $Z$ 后，常规更新给出 $b_8(0)=Y$。用 $Z$ 的补值替换它，四个答案都将翻转。在送达之前，对每个固定的 $K$，窗口 $(-1,0,1)$ 的记录两两相同。这一检验由 \nolinkurl{observer/local_return.py} 完成。

自主实现采用一个长度为 61 的环，以及一个半径为 1、空间齐次且不随时间变化的局部转移规则。每个元胞保存主层比特、相位计数器、固定角色和一个消息比特。初始化时，位置 3 被指定为观察者，位置 0 被指定为接收方。相位控制着消息的捕获、局部转发和写入。初始化后无需外部指令。对每一对世界，时刻 4 至 7 主层都相同，而完整的扩展状态保持不同。\nolinkurl{observer/autonomous.py} 的检验在时刻 8 得到 4/4 个正确输出。尚待解决的问题 R110.11 是在单一的、严格二元且齐次的规则~110 中实现同一套完整协议；上述构造使用的是扩展的有限字母表。对于孤立的确定性系统，一旦两段历史中的系统全状态完全相同，只有引入新的、能区分两者的输入，差异才可能再次出现。

\paragraph{形式化验证状态。}
在历史版本中，一个独立版本的观察模块在 Lean~4.33.1 中给出了九项 \texttt{\#print axioms} 报告，其中包括基于纤维的可预测性判据以及覆盖不充分的反例。在本专著的构建过程中，Lean~4.35.0 检查了导入 \texttt{Rule110Exact} 和 \texttt{Rule110FrontGlue} 的规范源码，并输出了所要求的全部 11/11 份公理列表。日志见附录~\ref{app:lean}。本节各有限实例另有相应的演化轨迹和穷举结果作为见证。
